%----------------------------------------------------------------------------------------
%   PACKAGES AND OTHER DOCUMENT CONFIGURATIONS
%----------------------------------------------------------------------------------------

\documentclass[bigger]{beamer}

\graphicspath{{./}} % Specifies where to look for included images

\usepackage{booktabs} % Allows the use of \toprule, \midrule and \bottomrule for better rules in tables

\mode<presentation>{}

%----------------------------------------------------------------------------------------
%   LAYOUT THEME
%----------------------------------------------------------------------------------------

\usetheme{Madrid}

%----------------------------------------------------------------------------------------
%   FONT THEME & FONTS
%----------------------------------------------------------------------------------------

\usefonttheme{default} % Typeset using the default sans serif font

%------------------------------------------------

\usepackage{palatino} % Use the Palatino font for serif text
\usepackage[default]{FiraSans} % Use the Fira Sans font for sans serif text

%----------------------------------------------------------------------------------------
%   INNER THEME
%----------------------------------------------------------------------------------------

\useinnertheme{circles}

%----------------------------------------------------------------------------------------
%   OUTER THEME
%----------------------------------------------------------------------------------------

\setbeamertemplate{footline}[page number] % Replace the footer line in all slides with a simple slide count
\setbeamertemplate{navigation symbols}{} % Remove the navigation symbols from the bottom of all slides

%----------------------------------------------------------------------------------------
%   PRESENTATION INFORMATION
%----------------------------------------------------------------------------------------

\title{Network-Based Diversification of Stock Portfolios}

\subtitle{Network Science: Project Presentation}

\author{João Gouveia, Gonçalo Rua}

\institute{Instituto Superior Técnico}

\date{\today}

%----------------------------------------------------------------------------------------

\begin{document}

%----------------------------------------------------------------------------------------
%   TITLE SLIDE
%----------------------------------------------------------------------------------------

\begin{frame}
    \titlepage
\end{frame}

%----------------------------------------------------------------------------------------
%   TABLE OF CONTENTS SLIDE
%----------------------------------------------------------------------------------------

\begin{frame}{Presentation Overview}
    \tableofcontents
\end{frame}

%----------------------------------------------------------------------------------------
%   PRESENTATION BODY SLIDES
%----------------------------------------------------------------------------------------

\section{Introduction}

\subsection{Motivation and previous work}

\begin{frame}{Introduction}{Motivation}
    DeFi contracts are vulnerable to attacks:
    \begin{itemize}
        \item<2-> Billions have been stolen.
        \item<3-> Exploits are mainly financial, not software-based.
        \item<4-> Leverage complex financial interactions among contracts, whose composition is poorly understood.
        \item<5-> Traditional bug-finding tools are not suitable to find these exploits.
    \end{itemize}

    \begin{exampleblock}{DeFi (Decentralized Finance)}
      Financial instruments built on top of blockchains. DeFi contracts make
      possible new financial instruments, that cannot be replicated with real
      world assets.
    \end{exampleblock}
\end{frame}

\begin{frame}{Introduction}{Previous Work}
    Previous work on formal verification of economic security in DeFi contracts:
    \begin{itemize}
        \item<2-> Focused on identifying known attacks.
        \item<3-> Lack the generality to discover new attacks.
    \end{itemize}

    \begin{exampleblock}{Formal Verification}
      Usage of mathematical models and proofs to prove that a system respects a
      given property.
    \end{exampleblock}
\end{frame}

\subsection{Contributions}

\begin{frame}{Introduction}{Contributions}
    This paper provides two main contributions:
    \begin{itemize}
        \item<2-> \textbf{EV} (Extractable Value): A formal definition for the economic security of composed contracts.
        \item<3-> \textbf{CFF} (Clockwork Finance Framework): Framework that relies on EV to prove bounds on the economic security of contract compositions.
    \end{itemize}
\end{frame}

%------------------------------------------------

\section{Extractable Value}

\begin{frame}{Extractable Value}{Informal Definition}
  Given any valid sequence of blocks, what is the maximum value that can be
  extracted by a specific participant in the system?
\end{frame}

\begin{frame}{Extractable Value}{Formal Definition}
    \begin{exampleblock}{(Simplified) Formal Definition}
        $\text{EV}(P,B,s) = \max_{(B_1, \ldots, B_k) \in B} \{ \sum_{a \in A_p} \text{balance}_k(a) - \text{balance}_0(a)\}$
    \end{exampleblock}
    \bigskip
    \begin{itemize}
        \item $P$: Specific participant in the system. Also referred to as player.
        \item $B$: Set of all valid block sequences of length $k$.
        \item $s$: Initial state, mapping of balances and data to accounts.
        \item $A_P$: Set of accounts belonging to participant $P$.
        \item $\text{balance}_{k}(a)$: Balance of account $a$ in state $s_k$.
    \end{itemize}
\end{frame}

%------------------------------------------------

\section{Clockwork Finance Framework}

\subsection{Use Cases}

\begin{frame}{Clockwork Finance Framework}{Use Cases}
    \begin{itemize}
        \item<1-> \textbf{DeFi contract designers}: Reason about the economic security of interactions between their contracts and other contracts.
        \item<2-> \textbf{Bots and miners}: Find profitable strategies in real-time.
        \item<3-> \textbf{Users}: Reason about the guarantees provided by the transactions they execute in the network, and the value at risk by exploits.
    \end{itemize}
\end{frame}

\subsection{Proving Economic Security}

\begin{frame}{Clockwork Finance Framework}{Proving Economic Security}
    \begin{itemize}
        \item<1-> CFF proves the economic security of a contract by proving its composability with the other contracts in the system.
        \item<2-> Contracts are modeled.
        \item<3-> Contract models fed to a solver.
        \item<4-> Solver either proves the safety of the contract, or returns an example showing why it is not safe.
    \end{itemize}

    \begin{exampleblock}{Composability}
      A contract $C_1$ is composable with $C_2$ if adding $C_2$ to the system
      does not give an adversary a significantly higher extractable value.
    \end{exampleblock}
\end{frame}

\begin{frame}{Clockwork Finance Framework}{Proving Economic Security}
    \begin{figure}
        \includegraphics[width=0.53\linewidth]{architecture.png}
        \caption{CFF's (simplified) architecture.}
    \end{figure}
\end{frame}

\subsection{Properties}

\begin{frame}{Clockwork Finance Framework}{Functional Properties}
    \begin{itemize}
        \item<1-> \textbf{Contract Completness}.
        \begin{itemize}
            \item<2-> Can model any DeFi contract with equivalent execution complexity to the native platform.
        \end{itemize}
    \end{itemize}
\end{frame}

\begin{frame}{Clockwork Finance Framework}{Functional Properties}
    \begin{itemize}
        \item<1-> \textbf{Attack Exhaustive by Construction}.
        \begin{itemize}
            \item<2-> Reasons about the all possible executions.
            \item<3-> Does not require manual coding of adversarial strategies.
            \item<4-> Can identify any attack expressible by its definitions.
            \item<5-> Yields false positives, not false negatives.
        \end{itemize}
    \end{itemize}
\end{frame}

\begin{frame}{Clockwork Finance Framework}{Usability Properties}
    \begin{itemize}
        \item<1-> \textbf{Modularity}: Models for a contract can be used by any CFF execution involving that contract.
        \item<2-> \textbf{Human Readability}: Models are typically easier to read than the corresponding contract source code.
    \end{itemize}
\end{frame}

%------------------------------------------------

\section{Review}

\subsection{The Good}

\begin{frame}{Review}{The Good}
    \begin{itemize}
        \item<1-> Introduces the notion of extractable value, and a new definition of contract composability.
        \item<2-> Both useful going forward, when reasoning about economic security.
    \end{itemize}
\end{frame}

\subsection{The Bad}

\begin{frame}{Review}{The Bad}
    \begin{itemize}
        \item<1-> The authors claim that CFF is useful to DeFi contract designers.
        \item<2-> The composability of a contract can be affected by the design flaws of other contracts.
        \begin{itemize}
            \item<3-> Namely by contracts added later to the system.
        \end{itemize}
        \item<4-> In practice, CFF does not seem to offer any useful guarantees to contract developers.
        \item<5-> Our contract might be provably economicaly safe today.
        \begin{itemize}
            \item<6-> Tommorow a new contract might be added to the system that completely negates this.
        \end{itemize}
    \end{itemize}
\end{frame}

\subsection{The Ugly}

\begin{frame}{Review}{The Ugly}
    \begin{itemize}
        \item<1-> CFF comes with 14 pre-modeled contracts.
        \item<2-> Any other contract must be manually modeled by the user.
        \item<3-> This causes two fundamental problems.
    \end{itemize}
\end{frame}

\begin{frame}{Review}{The Ugly}
    \textbf{Barrier of entry}.
    \begin{itemize}
        \item<2-> Users must learn how to model contracts before actually being able to use the tool.
        \item<3-> Not all users will be willing to do this.
        \item<4-> Widespread adoption becomes unlikely.
    \end{itemize}
\end{frame}

\begin{frame}{Review}{The Ugly}
    \textbf{Rotten foundation}.
    \begin{itemize}
        \item<2-> Manual modelling of the contracts introduces the possibility of making a wrong model.
        \item<3-> If the model itself is wrong, the proof is worthless.
        \item<4-> The danger of a wrong model is heightened by the fact that the target users of this tool will likely be new to formal verification.
    \end{itemize}
\end{frame}

%------------------------------------------------

\section{Questions}

\begin{frame}{Questions}
    \begin{itemize}
        \item Considering the issues with composability, will it ever be
        possible to develop a framework that provides relevant guarantees to
        contract developers?
        \item How can we solve the problems that come with manually modelling contracts?
    \end{itemize}
\end{frame}

\end{document} 
